% Generated draft: requires mathematical and source review.
\documentclass[11pt]{article}
\usepackage[margin=25mm]{geometry}
\usepackage{fontspec}
\setmainfont{Latin Modern Roman}
\usepackage{amsmath,amssymb,mathtools,unicode-math}
\setmathfont{Latin Modern Math}
\usepackage{enumitem,needspace,xurl,hyperref}
\setlength{\parindent}{0pt}
\setlength{\parskip}{5pt}
\setlength{\emergencystretch}{3em}
\newcommand{\N}{\mathbb{N}}
\newcommand{\Z}{\mathbb{Z}}
\newcommand{\Q}{\mathbb{Q}}
\newcommand{\R}{\mathbb{R}}
\newcommand{\C}{\mathbb{C}}
\newcommand{\F}{\mathbb{F}}
\newcommand{\A}{\mathbb{A}}
\newcommand{\PP}{\mathbb{P}}
\newcommand{\E}{\mathbb{E}}
\newcommand{\eps}{\varepsilon}
\newcommand{\dd}{\,\mathrm{d}}
\begin{document}
\section*{UnsolvedMath ID: 30002014}
\textbf{Source problem code:} catalogue entry
\subsection{Definitions and mathematical statement}
Let \(n\geq 1\) be an integer, let \(F\) be an arbitrary field, and let \(S_n\) denote the symmetric group on \(\{1,\ldots,n\}\). A partition \(\lambda\vdash n\) is a finite weakly decreasing sequence of positive integers with sum \(n\). Its Young diagram has \(\lambda_i\) boxes in row \(i\). A \(\lambda\)-tableau is a bijective filling of this diagram with \(1,\ldots,n\); its row tabloid \(\{t\}\) records the sets of entries in its respective rows. Let \(M_F^\lambda\) be the permutation \(F S_n\)-module with these tabloids as an \(F\)-basis. For a tableau \(t\), let \(C_t\leq S_n\) be the subgroup preserving each column's set of entries, and define its polytabloid by
\[
e_t=\sum_{\sigma\in C_t}\operatorname{sgn}(\sigma)\{\sigma t\}\in M_F^\lambda,
\]
where the signs are interpreted in \(F\). The Specht module \(S_F^\lambda\) is the \(F S_n\)-submodule spanned by all such \(e_t\).

\textbf{Homomorphism-space problem.} For every integer \(n\geq1\), every field \(F\), and every pair of partitions \(\lambda,\mu\vdash n\), determine the \(F\)-vector space
\[
\operatorname{Hom}_{F S_n}(S_F^\lambda,S_F^\mu)
=
\bigl\{f\in\operatorname{Hom}_F(S_F^\lambda,S_F^\mu):
 f(gv)=g f(v)\text{ for every }g\in S_n,\ v\in S_F^\lambda\bigr\}.
\]
A concrete determination would describe its dimension and a basis, or an equivalent explicit description of all its maps, uniformly in the stated data. No restriction to algebraically closed fields or to a particular characteristic is supplied.

\textbf{Associated decomposition-number target and unresolved qualification.} For an isomorphism class of simple \(F S_n\)-modules represented by \(D\), write
\[
d_{\lambda,D}=[S_F^\lambda:D]
\]
for the multiplicity of \(D\) among the composition factors of \(S_F^\lambda\). These multiplicities give a precise interpretation of the associated decomposition-number target. The supplied statement asks to compute the homomorphism spaces, ``thereby determining'' the associated decomposition-number problem. It supplies no theorem or procedure establishing that implication, and does not specify whether additional information or hypotheses are intended. Consequently the homomorphism-space task and the associated multiplicity target are recorded here without asserting that the former determines the latter. The source also does not specify the desired form of a general computation; the concrete output description above is an interpretation requiring review. The displayed range uses positive \(n\); the source does not explicitly address \(n=0\).
\subsection{Short English statement}
Determine all module homomorphisms between any two Specht modules for the same symmetric group over every field. The supplied record also asks for the associated decomposition numbers, interpreted as simple composition-factor multiplicities, but does not justify how the requested homomorphism spaces determine them. Its open-status label and literature assessment are supplied curation requiring independent review.
\subsection{Sources}
\begin{itemize}
\item Supplied local UnsolvedMath registry record, permanent ID \(30002014\), title \emph{Homomorphisms between Specht Modules}, source problem code \texttt{catalogue entry}, in the Oberwolfach Reports Open Problems collection. The supplied record gives proposed year 2012 and names no individual proposer.
\item \emph{Representations of Finite Groups} (2012), Oberwolfach Reports, pp.\ 963--1019, DOI \texttt{10.4171/owr/2012/16}. Supplied report reference: \url{https://doi.org/10.4171/owr/2012/16}. The record describes this collection as extracted from workshop open-problem discussions; no exact problem page or passage is supplied.
\item Supplied literature reference, labelled ``RoCK-block decomposition numbers'': \url{https://doi.org/10.1112/jlms.12852}.
\item Supplied literature reference, labelled ``Spin decomposition numbers'': \url{https://arxiv.org/abs/2303.00629}.
\item Literature assessment dated exactly \texttt{2026-08-21}: the supplied curation labels the problem open and states that general symmetric-group decomposition numbers remain outstanding, with substantial families and special blocks known but no uniform solution available. This is supplied dated curation derived from the OWR record, not an independently verified assessment. Its time-sensitive status and cited claims require independent verification.
\item The supplied statement has registry status \texttt{corrected\_verified}, accompanied by a claim of manual semantic, typographic, and LaTeX review and meaning-preserving cleanup. This records the supplied editorial qualification; it does not establish mathematical correctness or independent source verification here.
\end{itemize}
\end{document}
